\documentclass[12pt]{article}
\usepackage[utf8]{inputenc}
\usepackage[T1]{fontenc}
\usepackage{hyperref}
\title{Yaqui Language Resources for K-12 Curriculum Development}
\author{Compiled by Claude Sonnet 5.5 (claude-sonnet-5-5) for Andr\'as Kornai}
\date{}
\begin{document}
\maketitle

\section{Introduction}
Yaqui (Yoeme, Hiaki; ISO 639-3 yaq; Uto-Aztecan) is spoken in Sonora, Mexico, and Arizona, USA. This report lists only resources that appeared in search results. The Arizona variety was described in an NSF abstract as moribund.

\section{Primary Educational Resources}
\begin{itemize}
\item NSF project BCS-0843898, which planned a Hiaki teaching grammar with texts and exercises. I could not confirm that it was completed. \url{https://grantome.com/grant/NSF/BCS-0843898}
\end{itemize}

\section{Linguistic and Reference Materials}
\begin{itemize}
\item UCLA Yoeme (Yaqui) language site, whose database is described as drawn from the Hippocrene dictionary (Valenzuela, Molina, Shaul). \url{https://wil.cdh.ucla.edu/yaqui/?p=119}
\end{itemize}
The search summary also mentions \emph{Diccionario yaqui-espa\~nol y textos} (Universidad de Sonora and Plaza y Valdés, 2004) and Molina et al. (1999). I did not see their URLs, so there are no links.

\section{Religious and Community Resources}
None found in these searches. The summary mentions the Pascua Yaqui Tribe's language program only in passing, with no URL.

\section{Digital and Audio Resources}
\begin{itemize}
\item Endangered Languages Project entry for Yaqui. \url{https://endangeredlanguages.com/bn/node/84788}
\end{itemize}

\section{Curriculum Development Context}
Resources found are mostly academic and dictionary-oriented. I found no school textbooks, primers or programmes. The language is endangered and the Arizona and Sonora varieties differ in status. Any project should involve the Pascua Yaqui Tribe and Yaqui communities in Sonora.

\end{document}